\providecommand{\cO}{\mathcal{O}}

\section{The \texorpdfstring{$2A$}{2A} CHL Borcherds Constant}
\label{note:sec:direct-phi2A-chl-constant}

The order-two K3 input has three adjacent constants: the compact
Hodge supertrace of \(K3\times E\), the primitive CHL Jacobi
coefficient, and the Gritsenko-Clery diagonal-divisor catalogue
coefficient.  The BKM denominator uses the primitive CHL coefficient.
For the Nikulin involution the primitive CHL input is normalized by
\[
  \varphi^{\mathrm{CHL}}_{2}(\tau,z)\big|_{q^0}
  =
  \zeta^{-1} + c_2(0) + \zeta ,
  \qquad \zeta=e^{2\pi i z},
\]
and Borcherds' weight theorem reads this discriminant-zero coefficient.

\begin{theorem}[Direct \texorpdfstring{$2A$}{2A} CHL coefficient]
\label{note:thm:direct-phi2A-chl-c2}
Let \(g_2\) be the Mathieu class \(2A\), realised as a Nikulin
symplectic involution of a K3 surface, with frame shape \(1^8 2^8\).
Let \(\Phi_2\) be the primitive CHL Borcherds denominator obtained from
the corresponding level-two weak Jacobi input.  Then
\[
  c_2(0)=8,\qquad
  \kappa_{\mathrm{BKM}}(\Phi_2)=\frac{c_2(0)}{2}=4.
\]
Equivalently,
\[
  \varphi^{\mathrm{CHL}}_{2}(\tau,z)\big|_{q^0}
  =
  \zeta^{-1}+8+\zeta .
\]
On the five CHL levels \(N\in\{1,2,3,4,6\}\), the same convention gives
\[
  \begin{array}{c|ccccc}
  N & 1 & 2 & 3 & 4 & 6\\
  \hline
  c_N(0) & 10 & 8 & 6 & 4 & 2\\
  \kappa_{\mathrm{BKM}}(\Phi_N)=c_N(0)/2
    & 5 & 4 & 3 & 2 & 1
  \end{array}.
\]
\end{theorem}

\begin{proof}
Borcherds' singular-theta product formula gives
\[
  \mathrm{wt}(\mathrm{Bor}(\varphi))=\frac{c(0)}{2}
\]
for a weight-zero Jacobi input \(\varphi\), where \(c(0)\) is the
discriminant-zero Fourier coefficient of the input.  Gritsenko-Nikulin
construct the level-two primitive denominator \(\Phi_2\) in the CHL
ladder as the weight-four paramodular form.  Therefore
\[
  c_2(0)=2\,\mathrm{wt}(\Phi_2)=2\cdot 4=8,
  \qquad
  \kappa_{\mathrm{BKM}}(\Phi_2)=\mathrm{wt}(\Phi_2)=4.
\]

The theta-component extraction gives the same value from the \(2A\)
input.  In the level-two CHL normalization, write \(T_N\) for the
frame-shape trace datum and \(A_N\) for the odd theta-component
constant in the index-one theta decomposition.  The \(1^8 2^8\)
frame shape gives \(T_2=16\), and the \(2A\) odd theta component gives
\(A_2=4\).  The CHL cusp identity is
\[
  T_N=c_N(0)+2A_N.
\]
Hence
\[
  c_2(0)=T_2-2A_2=16-8=8.
\]
The primitive polar coefficients are fixed to one by the denominator
normalization, so the \(q^0\)-head is
\(\zeta^{-1}+8+\zeta\).

At \(N=1\), Eichler-Zagier gives
\(\phi_{0,1}|_{q^0}=\zeta^{-1}+10+\zeta\), and Borcherds-Gritsenko
gives \(\kappa_{\mathrm{BKM}}(\Delta_5)=10/2=5\).  Gritsenko-Nikulin
give weight \(3\) at \(N=3\), weight \(2\) at \(N=4\), and weight
\(1\) at \(N=6\).  Borcherds' formula then produces exactly the table
in the theorem.
\end{proof}

\paragraph{Separation of invariants.}
For \(K3\times E\) the relevant invariants split as follows:
\[
  \kappa_{\mathrm{cat}}(K3\times E)
  =\chi(\cO_{K3})\chi(\cO_E)=2\cdot 0=0,
  \qquad
  \kappa_{\mathrm{ch}}(K3\times E)=0,
\]
\[
  \kappa_{\mathrm{ch}}^{\mathrm{Heis}}(K3\times E)=3,
  \qquad
  \kappa_{\mathrm{BKM}}(\mathfrak g_{\Delta_5})=5,
  \qquad
  \kappa_{\mathrm{fiber}}=24.
\]
The equality \(5=3+2\) at the unorbifolded K3 fibre is a numerical
coincidence between two different constructions.  The CHL denominator
uses the Borcherds weight formula
\(\kappa_{\mathrm{BKM}}(\Phi_N)=c_N(0)/2\).

\paragraph{Normalizations.}
The singly twined K3 genus, the primitive CHL denominator input, and
the Gritsenko-Clery simplest-divisor catalogue are different Jacobi or
paramodular objects.  The singly twined Jacobi-input convention has
central coefficient \(10\) at \(N=1\), coefficient \(4\) at \(N=2\),
and coefficient \(2\) at \(N=3,4,6\).  The Gritsenko-Clery catalogue
is triple-indexed by \((t,N;k)\) and carries its own constants.  The
primitive CHL row used here is
\[
  \begin{array}{c|ccccc}
  N & 1 & 2 & 3 & 4 & 6\\
  \hline
  c_N(0) & 10 & 8 & 6 & 4 & 2\\
  \kappa_{\mathrm{BKM}}(\Phi_N) & 5 & 4 & 3 & 2 & 1
  \end{array}.
\]

\paragraph{Output scope.}
This computation belongs to the Borcherds-denominator specialisation of
\(\Phi_3\).  For a Calabi-Yau threefold the native output of the
CY-to-chiral construction is \(E_1\)-chiral; braided \(E_2\)-structures
live on centres, doubles, or module enhancements after the required
hypotheses are supplied.  In the toric \(\mathbb C^3\) comparison, the
CoHA is the positive half \(Y^+\); the \(\mathcal W_{1+\infty}\) object
belongs to the double or centre completion.

\paragraph{Anchors.}
The coefficient row is the CHL-denominator scope of
Gritsenko-Nikulin \(1995\), Part II, Theorem \(2.1\), combined with
Borcherds \(1998\), Theorem \(13.3\).  In this repository the
Borcherds-weight theorem and the closed-form \(c_N(0)\) lemma record
the same coefficient row.  The diagonal Siegel engine and the CY-D
stratification tests supply executable witnesses.
